% ===================================================================
%  EN49 附加定理类环境（不改 preamble.tex）
%  由 main.tex 与 src/tools/compile_chunk.py 生成的 wrap 文件共同 \input，
%  保证「整书编译」与「单块编译」两条路径都可见。
%
%  为什么需要：
%   1) 原书的 Remark 1–15 是**独立编号序列**（另立计数器），不是与定义/定理
%      共用的「章.节.序号」。若用 preamble 里的 remark 环境（挂在 result 计数器上），
%      印出来会是「注 2.1.3」，还会把紧随其后的定义、引理等挤后一位，与原书不符。
%   2) 原书另有若干**无编号**定理（如 §1.6 的 Kuw20/GS22 镜像对称定理、
%      §2.3 的 Abouzaid–Ganatra 定理）。用有编号环境会平白占号，同样错位。
%
%  用法：
%   \begin{remarkx} ... \end{remarkx}               → 注 1、注 2、…
%   \begin{thmstar}（Abouzaid--Ganatra, [AG]）... \end{thmstar}  → 定理. （无号）
%   其余无编号环境：propstar / lemstar / corstar / defstar / conjstar / examstar
% ===================================================================
\makeatletter
\@ifundefined{c@remarkx}{%
  \newtheorem{remarkx}{注}%
}{}
\newtheorem*{thmstar}{定理}
\newtheorem*{propstar}{命题}
\newtheorem*{lemstar}{引理}
\newtheorem*{corstar}{推论}
\newtheorem*{defstar}{定义}
\newtheorem*{conjstar}{猜想}
\newtheorem*{examstar}{例}
\makeatother
